\section{The finished investigation as a short paper}\label{ch18:sec:paper}
An investigation is finished only when it has been written up so that another reader can check it without seeing your notes or your conversations with an assistant. For a project like this one, a short paper with four parts works well.
\begin{itemize}
\item \emph{Introduction.} State the question: what happens to the iteration of a contraction when every update is made with an error? Explain why the first guess, convergence to the fixed point, is false, using the constant bias of Section~\ref{ch18:sec:guess}. Then say which questions replace it: how far from $x_*$ the points can be, and when they still converge.
\item \emph{Theory.} State the setting precisely: the space, the contraction factor $q$, and the model of the update errors, including the requirement that every $y_n$ lies in $X$. Then give Theorem~\ref{thm:inexact} with its proof, and the example of Section~\ref{ch18:sec:sharp} showing that the factor $1/(1-q)$ cannot be lowered.
\item \emph{Experiments.} Present the exact computation of Section~\ref{ch18:sec:experiment}, with the program and enough detail that anyone can run it again and get the same output. A computation that others can repeat in this way is called \emph{reproducible}. The experiments illustrate the proved statements; they do not replace the proofs.
\item \emph{Conclusion.} Say clearly what the results do not cover. The approximate points must stay in $X$. The factor $q$ must be less than $1$: for $q=1$ the sum $1+q+\cdots+q^{n-1}$ equals $n$, so the bound grows without limit, and the starting error is not reduced at all. And bounded update errors that do not tend to zero need not die out: the points may stay a fixed distance away from $x_*$, or fail to converge altogether, as the examples of Section~\ref{ch18:sec:sharp} show.
\end{itemize}

A paper should also say what is new in it. For this course project, an honest statement is that the project reconstructs and analyzes a standard argument about how errors in the individual updates affect an iteration. It must not claim that the result is new merely because the student, or the assistant, had not met it before. As we noted at the start of the book, a result you discover for yourself may be correct and still be well known.

Before claiming that a result is new in research, compare its exact statement and assumptions with the published literature, especially the \emph{primary sources}: the papers and books in which results first appeared. When one source describes a result from another, look up the result itself rather than relying on the description. Only after that comparison can you decide what, if anything, is new. Keep in mind that a search, whether in a library catalogue or by an assistant, that finds nothing tells you about that search. It does not prove that the result is new.
